\noindent
Radio direction finding is the process of estimating the direction from which a
radio signal arrive. It is an important capability in wireless systems, where
locating a transmitter can support monitoring, signal tracking, and a better
understanding of the radio environment. This project investigates a practical
approach to direction-of-arrival estimation using software-defined radio (SDR)
and beamforming.

\medskip
\noindent
The system is developed in three phases. First, a signal in the 2.4~GHz band is
transmitted and its in-phase and quadrature (I/Q) samples are acquired using an
ADALM-Pluto SDR. These samples provide the complex signal data needed for
digital processing. Second, beamforming techniques are explored to combine
signals received at multiple antenna elements and estimate the direction of
the source, with the aim of following changes in its position. Most of the calculation are done with
the help of our telecommunications knowledge, Acquierd with Mr Tihion. Finally, an
antenna is designed for the system, considering the requirements of operation
at the target frequency.

\medskip
\noindent
Together, these phases form an end-to-end study, from signal transmission and
acquisition to spatial processing and antenna design. The report presents the
approach and results of each phase, and discusses the practical constraints
encountered when implementing a direction-finding system with SDR hardware.